
        \documentclass[fleqn]{article}      
        \usepackage{amsmath}
        \usepackage{fontspec}
        \usepackage{kotex} % 한국어 지원  

        \begin{document}      
        쉬움: \\\\  
1. 다음 중 적분의 정의는 무엇인가요?  
   a) 함수의 기울기를 구하는 것  
   b) 함수 아래의 면적을 구하는 것  
   c) 함수의 최대값을 찾는 것  
   d) 함수의 최소값을 찾는 것  

2. ∫x dx의 결과는 무엇인가요?  
   a) x^2 + C  
   b) 2x + C  
   c) x + C  
   d) x^2  

3. 다음 중 정적분의 기호는 무엇인가요?  
   a) ∂  
   b) ∫  
   c) ∇  
   d) ∑  

4. ∫(3x^2) dx의 결과는 무엇인가요?  
   a) x^3 + C  
   b) 3x^3 + C  
   c) x^3 + 3C  
   d) 3x^3 + 3C  

5. 적분은 어떤 상황에서 주로 사용되나요?  
   a) 최대값을 찾기 위해  
   b) 면적을 구하기 위해  
   c) 기울기를 구하기 위해  
   d) 함수의 변화율을 구하기 위해  

중간: \\\\  
1. ∫(2x + 1) dx의 결과는 무엇인가요?  
   a) x^2 + x + C  
   b) 2x^2 + x + C  
   c) 2x^2 + 2x + C  
   d) 2x + C  

2. 다음 함수의 정적분을 구하세요: ∫(4x^3) dx  
   (주관식)  

3. ∫(sin(x)) dx의 결과는 무엇인가요?  
   a) -cos(x) + C  
   b) cos(x) + C  
   c) sin(x) + C  
   d) -sin(x) + C  

4. ∫(e^x) dx의 결과는 무엇인가요?  
   (주관식)  

5. ∫(1/x) dx의 결과는 무엇인가요?  
   a) ln|x| + C  
   b) e^x + C  
   c) x + C  
   d) 1/x + C  

어려움: \\\\  
1. ∫(x^2 + 2x + 1) dx의 결과를 구하세요.  
   (주관식)  

2. ∫(0부터 1까지) (2x + 1) dx의 결과는 무엇인가요?  
   (주관식)  

3. ∫(cos(x)) dx의 결과는 무엇인가요?  
   a) sin(x) + C  
   b) -sin(x) + C  
   c) cos(x) + C  
   d) -cos(x) + C  

4. ∫(x^3) dx의 결과는 무엇인가요?  
   a) (1/4)x^4 + C  
   b) (1/3)x^4 + C  
   c) (1/2)x^4 + C  
   d) x^4 + C  

5. ∫(2x) dx의 결과를 구하세요.  
   (주관식)  

      
        \end{document} 
    